%10/05/2004 at 16:00 
%
%Cinem�tica do Ponto Material: 
%Movimento Rectil�neo
% 
%Written by TORDAR 

\documentclass[12pt]{article}  

\pagestyle{empty}

\usepackage{graphicx} 
\usepackage{pst-all} 

%Let's go: 

\begin{document} 
\null\vskip4cm
\begin{pspicture}(0,0)(1,0.5)
\psset{unit=10cm}
\psline[linewidth=1pt,linecolor=black,linestyle=dashed,dash=4mm 1mm]{-}(0,0.3)(1,0.3)
\psline[linewidth=1pt,linecolor=black,linestyle=dashed,dash=4mm 1mm]{-}(0.3,0.3)(0.3,0)
\psline[linewidth=1pt,linecolor=black,linestyle=dashed,dash=4mm 1mm]{-}(0.6,0.3)(0.6,0)
\psline[linewidth=1pt,linecolor=black,linestyle=dashed,dash=4mm 1mm]{-}(0.95,0.3)(0.95,0)
\psline[linewidth=3pt,linecolor=blue]{<->}(0,0.5)(0,0)(1,0)
\psline[linewidth=3pt,linecolor=black](0,0)(0.3,0.3)(0.6,0.3)
\psline[linewidth=3pt,linecolor=black](0.6,0.01)(0.95,0.01)
\rput{0}(1.07,0){\scalebox{1.2}{$t$ [s]}}
\rput{0}(0.07,0.55){\scalebox{1.2}{$a(t)$ [m/\sii]}}
\rput{0}(0.3,-0.05){\scalebox{1.2}{10}}
\rput{0}(0.6,-0.05){\scalebox{1.2}{20}}
\rput{0}(0.95,-0.05){\scalebox{1.2}{30}}
\rput{0}(-0.05,0.3){\scalebox{1.2}{10}}
\end{pspicture}

\end{document} 
